\documentclass[10pt,a4paper]{amsart}
\usepackage{amsmath,amssymb,amsthm,graphicx}
% For commutative diagrams we use Xy-pic
\usepackage[all]{xy}

% We use 2cell for 2-commutative diagrams.
\xyoption{2cell}
% Theorem environments.
%
\theoremstyle{plain}
\newtheorem{theorem}[subsection]{Theorem}
\newtheorem{proposition}[subsection]{Proposition}
\newtheorem{lemma}[subsection]{Lemma}

\theoremstyle{definition}
\newtheorem{definition}[subsection]{Definition}
\newtheorem{example}[subsection]{Example}
\newtheorem{exercise}[subsection]{Exercise}
\newtheorem{situation}[subsection]{Situation}

\theoremstyle{remark}
\newtheorem{remark}[subsection]{Remark}
\newtheorem{remarks}[subsection]{Remarks}

\numberwithin{equation}{subsection}

% Macros
%
\def\lim{\mathop{\mathrm{lim}}\nolimits}
\def\colim{\mathop{\mathrm{colim}}\nolimits}
\def\Spec{\mathop{\mathrm{Spec}}}
\def\Hom{\mathop{\mathrm{Hom}}\nolimits}
\def\Ext{\mathop{\mathrm{Ext}}\nolimits}
\def\SheafHom{\mathop{\mathcal{H}\!\mathit{om}}\nolimits}
\def\SheafExt{\mathop{\mathcal{E}\!\mathit{xt}}\nolimits}
\def\Sch{\mathit{Sch}}
\def\Mor{\mathop{\mathrm{Mor}}\nolimits}
\def\Ob{\mathop{\mathrm{Ob}}\nolimits}
\def\Sh{\mathop{\mathit{Sh}}\nolimits}
\def\NL{\mathop{N\!L}\nolimits}
\def\CH{\mathop{\mathrm{CH}}\nolimits}
\def\proetale{{pro\text{-}\acute{e}tale}}
\def\etale{{\acute{e}tale}}
\def\QCoh{\mathit{QCoh}}
\def\Ker{\mathop{\mathrm{Ker}}}
\def\Im{\mathop{\mathrm{Im}}}
\def\Coker{\mathop{\mathrm{Coker}}}
\def\Coim{\mathop{\mathrm{Coim}}}

% Boxtimes
%
\DeclareMathSymbol{\boxtimes}{\mathbin}{AMSa}{"02}

%
% Macros for moduli stacks/spaces
%
\def\QCohstack{\mathcal{QC}\!\mathit{oh}}
\def\Cohstack{\mathcal{C}\!\mathit{oh}}
\def\Spacesstack{\mathcal{S}\!\mathit{paces}}
\def\Quotfunctor{\mathrm{Quot}}
\def\Hilbfunctor{\mathrm{Hilb}}
\def\Curvesstack{\mathcal{C}\!\mathit{urves}}
\def\Polarizedstack{\mathcal{P}\!\mathit{olarized}}
\def\Complexesstack{\mathcal{C}\!\mathit{omplexes}}
% \Pic is the operator that assigns to X its picard group, usage \Pic(X)
% \Picardstack_{X/B} denotes the Picard stack of X over B
% \Picardfunctor_{X/B} denotes the Picard functor of X over B
\def\Pic{\mathop{\mathrm{Pic}}\nolimits}
\def\Picardstack{\mathcal{P}\!\mathit{ic}}
\def\Picardfunctor{\mathrm{Pic}}
\def\Deformationcategory{\mathcal{D}\!\mathit{ef}}
\newenvironment{reference}{}{}
\newenvironment{slogan}{}{}
\newenvironment{history}{}{}
\usepackage{cleveref}
\crefname{originalsection}{original section}{original sections}
\Crefname{originalsection}{Original section}{Original sections}
\makeatletter
\newlabel{situation-regular-model}{{0C61}{source}{}{sourcepointer-1}{}}
\newlabel{situation-regular-model@cref}{{[originalsection][0][]0C61}{[1][0][]source}{}{}{}}
\makeatother
\crefname{originallemma}{original lemma}{original lemmas}
\Crefname{originallemma}{Original lemma}{Original lemmas}
\makeatletter
\newlabel{resolve-lemma-contract-ample}{{0C2M}{source}{}{sourcepointer-2}{}}
\newlabel{resolve-lemma-contract-ample@cref}{{[originallemma][0][]0C2M}{[1][0][]source}{}{}{}}
\makeatother
\crefname{originallemma}{original lemma}{original lemmas}
\Crefname{originallemma}{Original lemma}{Original lemmas}
\makeatletter
\newlabel{divisors-lemma-blow-up-pullback-effective-Cartier}{{0809}{source}{}{sourcepointer-3}{}}
\newlabel{divisors-lemma-blow-up-pullback-effective-Cartier@cref}{{[originallemma][0][]0809}{[1][0][]source}{}{}{}}
\makeatother
\crefname{originallemma}{original lemma}{original lemmas}
\Crefname{originallemma}{Original lemma}{Original lemmas}
\makeatletter
\newlabel{divisors-lemma-effective-Cartier-divisor-is-a-sum}{{0BCP}{source}{}{sourcepointer-4}{}}
\newlabel{divisors-lemma-effective-Cartier-divisor-is-a-sum@cref}{{[originallemma][0][]0BCP}{[1][0][]source}{}{}{}}
\makeatother
\crefname{originallemma}{original lemma}{original lemmas}
\Crefname{originallemma}{Original lemma}{Original lemmas}
\makeatletter
\newlabel{divisors-lemma-pullback-effective-Cartier-divisors}{{0C4U}{source}{}{sourcepointer-5}{}}
\newlabel{divisors-lemma-pullback-effective-Cartier-divisors@cref}{{[originallemma][0][]0C4U}{[1][0][]source}{}{}{}}
\makeatother
\crefname{originallemma}{original lemma}{original lemmas}
\Crefname{originallemma}{Original lemma}{Original lemmas}
\makeatletter
\newlabel{varieties-lemma-degree-birational-pullback}{{0AYU}{source}{}{sourcepointer-6}{}}
\newlabel{varieties-lemma-degree-birational-pullback@cref}{{[originallemma][0][]0AYU}{[1][0][]source}{}{}{}}
\makeatother
\title{Contracting an exceptional curve of the first kind on a regular model preserves the numerical type of the remaining components}
\author{The Stacks Project Authors}
\date{stacks-project:models.tex:local-snapshot-20261012; GNU FDL 1.2 or later}
\begin{document}
\maketitle
\noindent\textit{Editorial note.} The counted unit is the Stacks Project Lemma lemma-blowdown-regular-model at source lines 5623-5639 together with its proof at source lines 5641-5688, a straight adjacent pairing in the chapter file. The chapter is the upstream file models.tex; the pool directory carries the longer name stacks-models and is not used as the chapter name anywhere in this card.\par\medskip
\noindent\textit{Editorial note.} This card has no context span; the situation named by the statement and the proof is declared as a pointer with the kind section, because the assembler's whitelist of kinds has no entry for a situation.\par\medskip
\noindent\textit{Editorial note.} The retained proof quotes five results of other chapters, in Resolution of Surfaces, Divisors and Varieties, and names the situation of the chapter. All are declared as non-load-bearing navigation pointers with their stable Stacks tags looked up in the shared tag index, and none of their statements is reproduced inside the span.\par\medskip
\noindent\textit{Editorial note.} The retained proof span starts at line 5641, the line of the opening marker read from the source, and contains exactly one proof environment.\par\medskip
\noindent\textit{Editorial note.} The statement contains two enumerated assertions and a displayed list of data; both are part of the counted unit and are retained inside the statement span.\par\medskip
\noindent\textit{Editorial note.} No citation occurs in the retained spans, so no bibliography entry is supplied.\par\medskip
\noindent\textit{Allowed standard prerequisites (editorial).}\begin{itemize}
\item The regular-model situation, contractions of exceptional curves and pullbacks of Cartier divisors. The situation fixes a discrete valuation ring, a smooth projective curve over its fraction field, a regular proper model and the components of the special fibre with their multiplicities; an exceptional curve of the first kind is a smooth rational curve of self-intersection minus one and it may be contracted, the contraction being projective and keeping the surface regular and projective. For an effective Cartier divisor on a regular surface the pullback under such a contraction is computed component by component, the blowup case being recorded in the Divisors chapter, and the degree of a divisor on a component is invariant under the contraction for components disjoint from the contracted one. The author's own lemmas of the chapter and the quoted results of the Resolution of Surfaces, Divisors and Varieties chapters are the allowed inputs.. The counted lemma is the computational core of the section on blowing down exceptional curves: it says that contracting a minus-one curve leaves the numerical type of the other components unchanged, which is what makes the contraction step of the classification of models work.
\end{itemize}\medskip
\ifdefined\hypertarget\hypertarget{sourcepointer-1}{}\fi
\noindent\textit{Source reference.} Original section 0C61: focused-seeds/AGP/sources/stacks-snapshot-20261011/tags.txt, source lines 15281--15281. Pointer to the situation named by the retained statement and by the proof, declared with the kind section because the assembler's whitelist of kinds has no entry for a situation. The number is the stable Stacks tag of the target, looked up in the shared tag index.\par\medskip
\ifdefined\hypertarget\hypertarget{sourcepointer-2}{}\fi
\noindent\textit{Source reference.} Original lemma 0C2M: focused-seeds/AGP/sources/stacks-snapshot-20261011/tags.txt, source lines 15162--15162. Allowed input of the Resolution of Surfaces chapter quoted by the retained proof to contract the exceptional curve keeping the surface regular and projective; the pointer is registered so the reference resolves and the cited statement is not reproduced inside the span. The number is the stable Stacks tag of the target, looked up in the shared tag index.\par\medskip
\ifdefined\hypertarget\hypertarget{sourcepointer-3}{}\fi
\noindent\textit{Source reference.} Original lemma 0809: focused-seeds/AGP/sources/stacks-snapshot-20261011/tags.txt, source lines 10174--10174. Allowed input of the Divisors chapter quoted by the retained proof for the pullback of an effective Cartier divisor under a blowing up; the pointer is registered so the reference resolves and the cited statement is not reproduced inside the span. The number is the stable Stacks tag of the target, looked up in the shared tag index.\par\medskip
\ifdefined\hypertarget\hypertarget{sourcepointer-4}{}\fi
\noindent\textit{Source reference.} Original lemma 0BCP: focused-seeds/AGP/sources/stacks-snapshot-20261011/tags.txt, source lines 14287--14287. Allowed input of the Divisors chapter quoted by the retained proof to write an effective Cartier divisor as a sum of its parts; the pointer is registered so the reference resolves and the cited statement is not reproduced inside the span. The number is the stable Stacks tag of the target, looked up in the shared tag index.\par\medskip
\ifdefined\hypertarget\hypertarget{sourcepointer-5}{}\fi
\noindent\textit{Source reference.} Original lemma 0C4U: focused-seeds/AGP/sources/stacks-snapshot-20261011/tags.txt, source lines 15239--15239. Allowed input of the Divisors chapter quoted by the retained proof for the behaviour of effective Cartier divisors under pullback; the pointer is registered so the reference resolves and the cited statement is not reproduced inside the span. The number is the stable Stacks tag of the target, looked up in the shared tag index.\par\medskip
\ifdefined\hypertarget\hypertarget{sourcepointer-6}{}\fi
\noindent\textit{Source reference.} Original lemma 0AYU: focused-seeds/AGP/sources/stacks-snapshot-20261011/tags.txt, source lines 13802--13802. Allowed input of the Varieties chapter quoted by the retained proof for the invariance of the degree under a birational pullback, which is what turns the pullback computation into the displayed formula for the new intersection numbers; the pointer is registered so the reference resolves and the cited statement is not reproduced inside the span. The number is the stable Stacks tag of the target, looked up in the shared tag index.\par\medskip
\begin{lemma}
\label{lemma-blowdown-regular-model}
In Situation \ref{situation-regular-model} assume that $C_n$ is
an exceptional curve of the first kind. Let $f : X \to X'$ be the
contraction of $C_n$. Let $C'_i = f(C_i)$. Write $X'_k = \sum m'_i C'_i$.
Then $X'$, $C'_i$, $i = 1, \ldots, n' = n - 1$, and $m'_i = m_i$
is as in Situation \ref{situation-regular-model} and we have
\begin{enumerate}
\item for $i, j < n$ we have
$(C'_i \cdot C'_j) =
(C_i \cdot C_j) - (C_i \cdot C_n) (C_j \cdot C_n) /(C_n \cdot C_n)$,
\item for $i < n$ if $C_i \cap C_n \not = \emptyset$, then there are
maps $\kappa_i \leftarrow \kappa'_i \rightarrow \kappa_n$.
\end{enumerate}
Here $\kappa_i = H^0(C_i, \mathcal{O}_{C_i})$ and
$\kappa'_i = H^0(C'_i, \mathcal{O}_{C'_i})$.
\end{lemma}

\begin{proof}
By Resolution of Surfaces, Lemma \ref{resolve-lemma-contract-ample}
we can contract $C_n$ by a morphism $f : X \to X'$ such that $X'$ is
regular and is projective over $R$. Thus we see that $X'$ is as in
Situation \ref{situation-regular-model}.
Let $x \in X'$ be the image of $C_n$.
Since $f$ defines an isomorphism $X \setminus C_n \to X' \setminus \{x\}$
it is clear that $m'_i = m_i$ for $i < n$.

\medskip\noindent
Part (2) of the lemma is immediately clear from
the existence of the morphisms $C_i \to C'_i$ and $C_n \to x \to C'_i$.

\medskip\noindent
By Divisors, Lemma \ref{divisors-lemma-blow-up-pullback-effective-Cartier}
the pullback $f^{-1}C'_i$ is defined. By
Divisors, Lemma \ref{divisors-lemma-effective-Cartier-divisor-is-a-sum}
we see that $f^{-1}C'_i = C_i + e_i C_n$ for some $e_i \geq 0$. Since
$\mathcal{O}_X(C_i + e_i C_n) = \mathcal{O}_X(f^{-1}C'_i) =
f^*\mathcal{O}_{X'}(C'_i)$
(Divisors, Lemma \ref{divisors-lemma-pullback-effective-Cartier-divisors})
and since the pullback of an invertible sheaf restricts to the
trivial invertible sheaf on $C_n$ we see that
$$
0 = \deg_{C_n}(\mathcal{O}_X(C_i + e_i C_n)) =
(C_i + e_i C_n \cdot C_n) = (C_i \cdot C_n) + e_i(C_n \cdot C_n)
$$
As $f_j = f|_{C_j} : C_j \to C_j$ is a
proper birational morphism of proper curves
over $k$, we see that $\deg_{C'_j}(\mathcal{O}_{X'}(C'_i)|_{C'_j})$ is the
same as $\deg_{C_j}(f_j^*\mathcal{O}_{X'}(C'_i)|_{C'_j})$
(Varieties, Lemma \ref{varieties-lemma-degree-birational-pullback}).
Looking at the commutative diagram
$$
\xymatrix{
C_j \ar[r] \ar[d]_{f_j} & X \ar[d]^f \\
C'_j \ar[r] & X'
}
$$
and using Divisors, Lemma
\ref{divisors-lemma-pullback-effective-Cartier-divisors}
we see that
$$
(C'_i \cdot C'_j) = \deg_{C'_j}(\mathcal{O}_{X'}(C'_i)|_{C'_j}) =
\deg_{C_j}(\mathcal{O}_X(C_i + e_i C_n)) = (C_i + e_i C_n \cdot C_j)
$$
Plugging in the formula for $e_i$ found above we see that (1) holds.
\end{proof}

\end{document}
